\bta{2002}

\begin{enumerate}
	%\renewcommand{\labelenumi}{\arabic{enumi}.}
	% A(\Alph) a(\alph) I(\Roman) i(\roman) 1(\arabic)
	%设定全局标号series=example	%引用全局变量resume=example
	%[topsep=-0.3em,parsep=-0.3em,itemsep=-0.3em,partopsep=-0.3em]
	%可使用leftmargin调整列表环境左边的空白长度 [leftmargin=0em]
	\item
	填空题(本题共 5 小题，每小题 3 分，满分 15 分)
	\begin{enumerate}
		%\renewcommand{\labelenumi}{\arabic{enumi}.}
		% A(\Alph) a(\alph) I(\Roman) i(\roman) 1(\arabic)
		%设定全局标号series=example	%引用全局变量resume=example
		%[topsep=-0.3em,parsep=-0.3em,itemsep=-0.3em,partopsep=-0.3em]
		%可使用leftmargin调整列表环境左边的空白长度 [leftmargin=0em]
		\item
		 $\lim\limits _{n \rightarrow \infty} \cos \frac{1}{2} \cos \frac{1}{4} \cdots \cos \frac{1}{2^{n}}=$ \tk{$\sin1$} 。
		
		
		\item
			 $\int \frac{\cos x d x}{1+e^{\sin x}}$ \tk{$\sin x-\ln\left(1+e^{\sin x}\right)+C$} 。
		
		
		\item
			设 $z=z(x, y)$, 且有 $y z+z x+x y=1$, 则 $d z=$
			\tk{$\displaystyle-\frac{y+z}{x+y}\,dx-\frac{x+z}{x+y}\,dy$} 。
		
		
		\item
			含参变量广义积分 $\int_{0}^{+\infty} \frac{\sin \frac{1}{x}}{x^{p}} d x$ 的收敛域为 \tk{$0<p<2$} 。
		
		
		\item
			微分方程 $y^{\prime \prime}+4 y^{\prime}+4 y=e^{x}$ 的通解为
			\tk{$\displaystyle y=(C_1+C_2x)e^{-2x}+\frac19e^x$} 。
		
		
		
		
	\end{enumerate}

\banswer{
	\begin{enumerate}
		%\renewcommand{\labelenumi}{\arabic{enumi}.}
		\item
		$\sin1$。反复使用二倍角公式 $\sin t=2\sin(t/2)\cos(t/2)$，
		可得
		\[
		\sin1
		=2^n\sin\frac1{2^n}
		\prod_{k=1}^{n}\cos\frac1{2^k},
		\]
		因此
		\[
		\prod_{k=1}^{n}\cos\frac1{2^k}
		=\frac{\sin1}{2^n\sin(1/2^n)}.
		\]
		由于 $2^n\sin(1/2^n)\to1$，所求极限为 $\sin1$。

		\item
		$\sin x-\ln\left(1+e^{\sin x}\right)+C$。令 $u=\sin x$，则
		$du=\cos x\,dx$，于是
		\[
		\int\frac{\cos x\,dx}{1+e^{\sin x}}
		=\int\frac{du}{1+e^u}.
		\]
		注意
		\[
		\frac{d}{du}\left[u-\ln(1+e^u)\right]
		=1-\frac{e^u}{1+e^u}
		=\frac1{1+e^u},
		\]
		回代 $u=\sin x$ 即得原积分。

		\item
		$dz=-\frac{y+z}{x+y}\,dx-\frac{x+z}{x+y}\,dy$。
		对隐式方程 $yz+zx+xy=1$ 取全微分：
		\[
		(z+y)\,dx+(z+x)\,dy+(x+y)\,dz=0,
		\]
		当 $x+y\ne0$ 时即解出上式。

		\item
		$0<p<2$。令 $t=1/x$，则 $x=1/t$、$dx=-dt/t^2$，从而
		\[
		\int_0^{+\infty}\frac{\sin(1/x)}{x^p}\,dx
		=\int_0^{+\infty}t^{p-2}\sin t\,dt.
		\]
		在 $t\to0^{+}$ 时 $\sin t\sim t$，被积函数等价于
		$t^{p-1}$，故要求 $p>0$；在 $t\to+\infty$ 时，
		由 Dirichlet 判别法，$t^{p-2}\sin t$ 的积分收敛需要
		$p-2<0$，即 $p<2$（$p=2$ 时 $\int^{+\infty}\sin t\,dt$ 发散）。
		所以收敛域为 $0<p<2$；进一步，$0<p<1$ 时绝对收敛，
		$1\leqslant p<2$ 时条件收敛。

		\item
		$y=(C_1+C_2x)e^{-2x}+\frac19e^x$。齐次方程的特征方程为
		\[
		r^2+4r+4=(r+2)^2=0,
		\]
		有二重根 $-2$，故 $y_h=(C_1+C_2x)e^{-2x}$。
		因 $1$ 不是特征根，对非齐次项设 $y^{*}=Ae^x$，代入得
		\[
		(A+4A+4A)e^x=e^x,
		\]
		所以 $A=\frac19$。
	\end{enumerate}
}

\item 	
选择题 (本题共 5 小题，每小题 3 分，满分 15 分)	
\begin{enumerate}
	%\renewcommand{\labelenumi}{\arabic{enumi}.}
	% A(\Alph) a(\alph) I(\Roman) i(\roman) 1(\arabic)
	%设定全局标号series=example	%引用全局变量resume=example
	%[topsep=-0.3em,parsep=-0.3em,itemsep=-0.3em,partopsep=-0.3em]
	%可使用leftmargin调整列表环境左边的空白长度 [leftmargin=0em]
	\item
	设 $\left\{\begin{array}{ll}a x+b & x \leqslant 0 \\ x^{2} \cos \frac{1}{x} & x>0\end{array}\right.$ 在 $x=0$ 处可导, 则 \xzanswer{B}

	
\fourchoices
{$a=1, b=0$}
{$a=0, b=0$}
{$a=1, b=1$}
{$a=0, b=1$}

	
\item 
	$f(x)$ 在区间 $[-L, L]$ 内具有二阶导数，且 $f^{\prime \prime}(x)>0, \lim\limits _{x \rightarrow 0} \frac{f(x)}{x}=1$, 则 \xzanswer{A}


\fourchoices
{在 $(-L, 0)$ 和 $0, L$ 内均有 $f(x)>x$}
{在 $(-L, 0)$ 和 $0, L$ 内均有 $f(x)<x$}
{在 $(-L, 0)$ 内, $f(x)>x$; 在 $0, L$ 内, $f(x)<x$}
{在 $(-L, 0)$ 内, $f(x)<x$; 在 $0, L$ 内, $f(x)>x$}



\item
	设 $L$ 为圆周
	$\left\{\begin{array}{l}
	x^{2}+y^{2}+z^{2}=a^{2},\\
	x+y+z=0,
	\end{array}\right.$
	则曲线积分 $\int_{L}\left(x^{4}+2 y^{2} z^{2}\right) d L=$ \xzanswer{B}


\fourchoices
{$\frac{\pi a^{5}}{3}$}
	{$\frac{2 \pi a^{5}}{3}$}
{$\pi a^{5}$}
{$2 \pi a^{5}$}



\item
	下列级数中, 绝对收敛的级数是 \xzanswer{C}


\fourchoices
{$\sum\limits_{n=1}^{\infty}\left(e^{\frac{1}{n}}-1\right)$}
{$\sum\limits_{n=1}^{\infty} \frac{\sin n}{n}$}
{$\sum\limits_{n=1}^{\infty} \frac{(-1)^{n}}{e^{\sqrt{n}}+1}$}
{$\sum\limits_{n=1}^{\infty} \frac{(-1)^{n}}{\sqrt[n]{e}+1}$}


\item 
	$ a+|x|=\pi-\frac{4}{\pi} \sum\limits_{n=1}^{\infty} \frac{\cos (2 n-1) x}{(2 n-1)^{2}}$, 其中 $-\pi \leqslant x \leqslant \pi, a$ 为常数， 则 $a=$ \xzanswer{A}


\fourchoices
{$\frac{\pi}{2}$}
{$-\frac{\pi}{2}$}
{$\pi$}
{$-\pi$}



	
\end{enumerate}

\banswer{
	\begin{enumerate}
		%\renewcommand{\labelenumi}{\arabic{enumi}.}
		\item
		\textbf{B}。记分段函数为 $\varphi(x)$。可导必先连续，因此
		\[
		b=\varphi(0)=\lim_{x\to0^{+}}x^2\cos\frac1x=0.
		\]
		左导数为 $a$，而右导数为
		\[
		\lim_{x\to0^{+}}
		\frac{x^2\cos(1/x)-0}{x}
		=\lim_{x\to0^{+}}x\cos\frac1x=0,
		\]
		左右导数相等要求 $a=0$。故 $a=0$、$b=0$。

		\item
		\textbf{A}。$f$ 在 $0$ 处二阶可导，因而连续。由
		$\lim\limits_{x\to0}f(x)/x=1$ 先得 $f(0)=0$，再得
		\[
		f'(0)=\lim_{x\to0}\frac{f(x)-f(0)}x=1.
		\]
		由于 $f''(x)>0$，$f$ 严格凸，其图像除切点外严格位于
		任一切线之上。而 $x=0$ 处的切线为
		\[
		y=f(0)+f'(0)x=x,
		\]
		所以对 $x\in(-L,0)\cup(0,L)$ 均有 $f(x)>x$。

		\item
		\textbf{B}。$L$ 是半径为 $a$ 的大圆，弧长 $2\pi a$。在 $L$ 上由
		$x+y+z=0$ 得
		\[
		0=(x+y+z)^{2}=a^{2}+2(xy+yz+zx),
		\]
		即 $xy+yz+zx=-\frac{a^{2}}{2}$。于是
		\begin{align*}
		x^{2}y^{2}+y^{2}z^{2}+z^{2}x^{2}
		&=(xy+yz+zx)^{2}-2xyz(x+y+z)=\frac{a^{4}}{4},\\
		x^{4}+y^{4}+z^{4}
		&=\left(x^{2}+y^{2}+z^{2}\right)^{2}
		-2\left(x^{2}y^{2}+y^{2}z^{2}+z^{2}x^{2}\right)\\
		&=a^{4}-\frac{a^{4}}{2}=\frac{a^{4}}{2}.
		\end{align*}
		两式右端在 $L$ 上都是常数，而 $x,y,z$ 的地位完全对称，故
		\[
		\int_Lx^4\,dL=\frac{a^{4}}{6}\cdot2\pi a=\frac{\pi a^5}{3},\qquad
		\int_Ly^2z^2\,dL=\frac{a^{4}}{12}\cdot2\pi a=\frac{\pi a^5}{6},
		\]
		因此 $\int_L\left(x^4+2y^2z^2\right)dL=\frac{2\pi a^5}{3}$
		（已用参数化数值积分核对）。

		\item
		\textbf{C}。逐项判断：由 $e^{1/n}-1\sim\frac1n$ 知 A 发散；
		B 由 Dirichlet 判别法收敛，但
		\[
		\sum\frac{|\sin n|}{n}
		\geqslant\sum\frac{\sin^2n}{n},
		\]
		且后者利用 $\sin^2n=(1-\cos2n)/2$ 可知发散，所以
		B 不绝对收敛。对 C，
		\[
		0<\frac1{e^{\sqrt n}+1}<e^{-\sqrt n},
		\]
		而 $\sum e^{-\sqrt n}$ 收敛，故 C 绝对收敛。D 的通项
		绝对值趋于 $\frac12$，不满足通项趋零的必要条件。

		\item
		\textbf{A}。在 $[-\pi,\pi]$ 上对等式两边积分，所有余弦项的积分
		均为零，因此
		\[
		2\pi a+\int_{-\pi}^{\pi}|x|\,dx
		=\int_{-\pi}^{\pi}\pi\,dx,
		\]
		即 $2\pi a+\pi^2=2\pi^2$，从而 $a=\frac\pi2$。
	\end{enumerate}
}


\newpage
\item 
(3 小题，每小题 6 分，共 18 分)
\begin{enumerate}
	%\renewcommand{\labelenumi}{\arabic{enumi}.}
	% A(\Alph) a(\alph) I(\Roman) i(\roman) 1(\arabic)
	%设定全局标号series=example	%引用全局变量resume=example
	%[topsep=-0.3em,parsep=-0.3em,itemsep=-0.3em,partopsep=-0.3em]
	%可使用leftmargin调整列表环境左边的空白长度 [leftmargin=0em]
	\item
计算定积分 $\int_{-\frac{\pi}{2}}^{\frac{\pi}{2}} \frac{\cos x}{1+e^{x}} d x$。
	
	
	\item
	计算积分 $\int_{0}^{1}(1-\sqrt{x})^{n} d x$。
	
	
	\item
	设函数 $f(x)$ 满足 $f^{\prime \prime}(x)+\left[f^{\prime}(x)\right]^{2}=\sin x$, 且 $f^{\prime}(0)=0$ 。证明: $x=0$ 是 $f(x)$ 的拐点。
	
	
	
	
\end{enumerate}


\banswer{
	\begin{enumerate}
		\item
		记
		\[
		h(x)=\frac{\cos x}{1+e^x}.
		\]
		把 $h(-x)$ 的分母同乘 $e^x$，有
		\[
		h(x)+h(-x)
		=\cos x\left(\frac1{1+e^x}+\frac1{1+e^{-x}}\right)
		=\cos x\left(\frac1{1+e^x}+\frac{e^x}{1+e^x}\right)
		=\cos x,
		\]
		又因为积分区间关于原点对称，
		\[
		\begin{aligned}
		2\int_{-\pi/2}^{\pi/2}h(x)\,dx
		&=\int_{-\pi/2}^{\pi/2}\bigl[h(x)+h(-x)\bigr]\,dx\\
		&=\int_{-\pi/2}^{\pi/2}\cos x\,dx=2.
		\end{aligned}
		\]
		所以
		\[
		\int_{-\pi/2}^{\pi/2}\frac{\cos x}{1+e^x}\,dx
		=\boxed{1}.
		\]
		\item
		令 $t=\sqrt{x}$，则 $x=t^2$、$dx=2t\,dt$，所以
		\[
		\int_0^1(1-\sqrt{x})^n\,dx
		=2\int_0^1t(1-t)^n\,dt.
		\]
		对后一积分分部积分：
		\[
		\begin{aligned}
		\int_0^1t(1-t)^n\,dt
		&=\left[-\frac{t(1-t)^{n+1}}{n+1}\right]_0^1
		+\frac1{n+1}\int_0^1(1-t)^{n+1}\,dt\\
		&=\frac1{(n+1)(n+2)}.
		\end{aligned}
		\]
		故
		\[
		\boxed{\displaystyle
		\int_0^1(1-\sqrt{x})^n\,dx
		=\frac{2}{(n+1)(n+2)}}.
		\]
		\item
		先在原方程中令 $x=0$。由 $f'(0)=0$ 得
		\[
		f''(0)+[f'(0)]^2=0,
		\]
		即 $f''(0)=0$。又由
		\[
		f''(x)=\sin x-[f'(x)]^2
		\]
		可知 $f''$ 在 $0$ 附近可微。对方程求导：
		\[
		f'''(x)+2f'(x)f''(x)=\cos x,
		\]
		令 $x=0$ 得 $f'''(0)=1$。于是
		\[
		\lim_{x\to0}\frac{f''(x)}x
		=\lim_{x\to0}\frac{f''(x)-f''(0)}{x-0}
		=f'''(0)=1.
		\]
		因此存在 $\delta>0$，使得当 $-\delta<x<0$ 时
		$f''(x)<0$，当 $0<x<\delta$ 时 $f''(x)>0$。
		曲线在 $x=0$ 两侧的凹凸性发生改变，故
		\[
		\boxed{x=0\text{ 是 }f(x)\text{ 的拐点}}.
		\]
	\end{enumerate}
}



\item 
(4 小题，每小题 7 分，共 28 分)
\begin{enumerate}
	%\renewcommand{\labelenumi}{\arabic{enumi}.}
	% A(\Alph) a(\alph) I(\Roman) i(\roman) 1(\arabic)
	%设定全局标号series=example	%引用全局变量resume=example
	%[topsep=-0.3em,parsep=-0.3em,itemsep=-0.3em,partopsep=-0.3em]
	%可使用leftmargin调整列表环境左边的空白长度 [leftmargin=0em]
	\item
计算极限 $\lim\limits _{n \rightarrow \infty} \sqrt[n]{1+e^{n x}+e^{-n x}}$。
	
	
	\item
	计算曲线积分 $I=\oint_{L} z^{2} d x+\left(x^{2}+x y-x\right) d y+2 x z d z$, 其中 $L$ 是抛物面 $z=x^{2}+y^{2}$ 与圆柱面 $x^{2}+4 y^{2}=1$ 的交线。从 $z$ 轴正方向向下看, $L$ 为顺时针方向。
	
	
	\item
	把 $y=\arctan \frac{3+x}{3-x}$ 展为 $x$ 的幂级数, 并求收敛域。
	
	
	\item
	求微分方程 $\left(x-x^{3} y^{2} \ln y\right) y^{\prime}=2 y$ 的通解。
	
	
	
	
\end{enumerate}

\banswer{
	\begin{enumerate}
		\item
		分三种情况讨论。若 $x>0$，从根式内提取 $e^{nx}$：
		\[
		\sqrt[n]{1+e^{nx}+e^{-nx}}
		=e^x\left(1+e^{-nx}+e^{-2nx}\right)^{1/n}
		\longrightarrow e^x.
		\]
		若 $x<0$，提取 $e^{-nx}$，同理极限为 $e^{-x}$。
		若 $x=0$，原式为 $3^{1/n}\to1$。综上，
		\[
		\boxed{\displaystyle
		\lim_{n\to\infty}\sqrt[n]{1+e^{nx}+e^{-nx}}
		=e^{|x|}}.
		\]
		\item
		令
		\[
		\bm F=(P,Q,R)=(z^2,x^2+xy-x,2xz).
		\]
		直接计算旋度：
		\[
		\nabla\times\bm F
		=\left(R_y-Q_z,\ P_z-R_x,\ Q_x-P_y\right)
		=(0,0,2x+y-1).
		\]
		取抛物面 $z=x^2+y^2$ 位于椭圆柱内的部分为跨接曲面。
		题设从 $z$ 轴正方向向下看为顺时针，因此按右手法则应取
		向下法向量
		\[
		\bm n\,dS=(2x,2y,-1)\,dx\,dy.
		\]
		由 Stokes 公式，
		\[
		\begin{aligned}
		I
		&=\iint_{x^2+4y^2\leqslant1}
		(0,0,2x+y-1)\cdot(2x,2y,-1)\,dx\,dy\\
		&=\iint_{x^2+4y^2\leqslant1}(1-2x-y)\,dx\,dy.
		\end{aligned}
		\]
		积分区域关于 $x$ 轴和 $y$ 轴均对称，故 $x,y$ 的积分为零。
		椭圆的两个半轴为 $1$ 和 $1/2$，面积为 $\pi/2$，所以
		\[
		\boxed{I=\frac{\pi}{2}}.
		\]
		\item
		由正切加法公式，在 $|x|<3$ 时
		\[
		\arctan\frac{3+x}{3-x}
		=\frac{\pi}{4}+\arctan\frac{x}{3}.
		\]
		再用
		\[
		\arctan u=\sum_{n=0}^{\infty}
		\frac{(-1)^nu^{2n+1}}{2n+1},\qquad |u|<1,
		\]
		得到
		\[
		\boxed{\displaystyle
		\arctan\frac{3+x}{3-x}
		=\frac{\pi}{4}+
		\sum_{n=0}^{\infty}
		\frac{(-1)^n x^{2n+1}}{(2n+1)3^{2n+1}}},
		\qquad |x|<3.
		\]
		由端点处的交错级数判别，幂级数在 $x=\pm3$ 均收敛，
		所以级数本身的收敛区间为 $[-3,3]$。在 $x=-3$ 时其和为
		$0$，与原函数相同；在 $x=3$ 时其和为 $\pi/2$，但原式
		分母为零。因此作为原函数的展开式，等式成立于
		\[
		\boxed{[-3,3)}.
		\]
		\item
		因为出现 $\ln y$，先在 $y>0$ 的区间上求解。把 $x$
		看成 $y$ 的函数：
		\[
		\frac{dx}{dy}
		=\frac{x-x^3y^2\ln y}{2y}
		=\frac{x}{2y}-\frac12x^3y\ln y.
		\]
		这是关于 $x(y)$ 的 Bernoulli 方程。令 $u=x^{-2}$，则
		\[
		u'=-2x^{-3}\frac{dx}{dy}
		=-\frac{u}{y}+y\ln y,
		\]
		即
		\[
		u'+\frac1yu=y\ln y.
		\]
		积分因子为 $\exp\!\int dy/y=y$，故
		\[
		(yu)'=y^2\ln y.
		\]
		分部积分得到
		\[
		yu=\int y^2\ln y\,dy
		=\frac{y^3}{3}\ln y-\frac{y^3}{9}+C.
		\]
		代回 $u=x^{-2}$，通解为
		\[
		\boxed{\displaystyle
		\frac{y}{x^2}-\frac{y^3}{3}\ln y+\frac{y^3}{9}=C},
		\qquad y>0,\quad x\ne0.
		\]
	\end{enumerate}
}


\item 
(3 小题，每小题 8 分，共 24 分)
\begin{enumerate}
	%\renewcommand{\labelenumi}{\arabic{enumi}.}
	% A(\Alph) a(\alph) I(\Roman) i(\roman) 1(\arabic)
	%设定全局标号series=example	%引用全局变量resume=example
	%[topsep=-0.3em,parsep=-0.3em,itemsep=-0.3em,partopsep=-0.3em]
	%可使用leftmargin调整列表环境左边的空白长度 [leftmargin=0em]
	\item
设曲面 $S: \sqrt{x}+\sqrt{y}+\sqrt{z}=\sqrt{a}(a>0)$, 在 $S$ 上求一切平面, 使此切面与三坐标 所围成的四面体体积最大，并求四面体体积的最大值。
	
\item
	设 $f(x)$ 是以 $2 \pi$ 为周期的函数, 在 $(-\pi, \pi]$ 内的表达式为
\[
f(x)=\left\{\begin{array}{ll}
	0 & -\pi<x<0 \\
	1 & 0 \leqslant x \leqslant 1 \\
	0 & 1<x \leqslant \pi
\end{array}\right.
\]	
将 $f(x)$ 展为傅里叶级数 (说明收敛情况), 并求 $\sum\limits_{n=1}^{\infty} \frac{\sin n}{n}$ 与 $\sum\limits_{n=1}^{\infty} \frac{1-\cos n}{n^{2}}$。 

\item
设区域 $\Omega$ 由曲面 $x=0, y=0, x+y=1, z(x+y)=1$ 及 $z=1$ 围成。
\begin{enumerate}
	%\renewcommand{\labelenumi}{\arabic{enumi}.}
	% A(\Alph) a(\alph) I(\Roman) i(\roman) 1(\arabic)
	%设定全局标号series=example	%引用全局变量resume=example
	%[topsep=-0.3em,parsep=-0.3em,itemsep=-0.3em,partopsep=-0.3em]
	%可使用leftmargin调整列表环境左边的空白长度 [leftmargin=0em]
	\item
求 $\Omega$ 的体积 $V$。
\item 
证明 $\iiint_{\Omega} \frac{d x d y d z}{x^{2}+y^{2}+z^{2}} \leqslant \frac{V}{2}$。
	
	
	
\end{enumerate}




	
\end{enumerate}

\banswer{
	\begin{enumerate}
		\item
		设切点为 $(x_0,y_0,z_0)$，并令
		\[
		p=\sqrt{x_0},\qquad q=\sqrt{y_0},\qquad r=\sqrt{z_0}.
		\]
		则 $p,q,r>0$ 且 $p+q+r=\sqrt a$。记
		\[
		F(x,y,z)=\sqrt x+\sqrt y+\sqrt z-\sqrt a.
		\]
		由
		\[
		\nabla F(x_0,y_0,z_0)
		=\left(\frac1{2p},\frac1{2q},\frac1{2r}\right),
		\]
		切平面方程为
		\[
		\frac{x-p^2}{2p}+\frac{y-q^2}{2q}
		+\frac{z-r^2}{2r}=0,
		\]
		即
		\[
		\frac{x}{p}+\frac{y}{q}+\frac{z}{r}=\sqrt a.
		\]
		该平面在三个坐标轴上的截距分别为
		$\sqrt a\,p,\sqrt a\,q,\sqrt a\,r$，因而所围四面体体积为
		\[
		V=\frac16(\sqrt a\,p)(\sqrt a\,q)(\sqrt a\,r)
		=\frac{a^{3/2}}6pqr.
		\]
		由算术--几何平均值不等式，
		\[
		pqr\leqslant
		\left(\frac{p+q+r}{3}\right)^3
		=\frac{a^{3/2}}{27},
		\]
		等号当且仅当 $p=q=r=\sqrt a/3$，即切点为
		$(a/9,a/9,a/9)$。此时所求切平面及最大体积为
		\[
		\boxed{x+y+z=\frac a3},\qquad
		\boxed{V_{\max}=\frac{a^3}{162}}.
		\]
		\item
		按 $2\pi$ 周期傅里叶系数的定义，
		\[
		\begin{aligned}
		a_0&=\frac1\pi\int_{-\pi}^{\pi}f(x)\,dx
		=\frac1\pi,\\
		a_n&=\frac1\pi\int_{-\pi}^{\pi}f(x)\cos(nx)\,dx
		=\frac1\pi\int_0^1\cos(nx)\,dx
		=\frac{\sin n}{\pi n},\\
		b_n&=\frac1\pi\int_{-\pi}^{\pi}f(x)\sin(nx)\,dx
		=\frac1\pi\int_0^1\sin(nx)\,dx
		=\frac{1-\cos n}{\pi n}.
		\end{aligned}
		\]
		因此傅里叶级数为
		\[
		\boxed{\displaystyle
		f(x)\sim\frac1{2\pi}
		+\frac1\pi\sum_{n=1}^{\infty}
		\left(\frac{\sin n}{n}\cos nx
		+\frac{1-\cos n}{n}\sin nx\right)}.
		\]
		按 Dirichlet 收敛定理，该级数在连续点收敛到 $f(x)$；
		在跳跃点 $x\equiv0,1\pmod{2\pi}$ 收敛到左右极限的平均值
		$1/2$。具体地，在 $0<x<1$ 时收敛到 $1$，在其余连续点
		收敛到 $0$。

		令 $x=0$。因为此处级数收敛到 $1/2$，
		\[
		\frac12=\frac1{2\pi}
		+\frac1\pi\sum_{n=1}^{\infty}\frac{\sin n}{n}.
		\]
		故
		\[
		\boxed{\displaystyle
		\sum_{n=1}^{\infty}\frac{\sin n}{n}
		=\frac{\pi-1}{2}}.
		\]
		另一方面，经典正弦级数
		\[
		\sum_{n=1}^{\infty}\frac{\sin(nt)}n
		=\frac{\pi-t}{2},\qquad 0<t<2\pi,
		\]
		从 $0$ 到 $1$ 积分（积分后的级数绝对收敛），得到
		\[
		\begin{aligned}
		\sum_{n=1}^{\infty}\frac{1-\cos n}{n^2}
		&=\int_0^1\frac{\pi-t}{2}\,dt\\
		&=\boxed{\frac{\pi}{2}-\frac14}.
		\end{aligned}
		\]
		\item
		作变量代换
		\[
		s=x+y,\qquad t=x.
		\]
		则 $x=t,\ y=s-t$，且
		\[
		\left|\frac{\partial(x,y)}{\partial(s,t)}\right|=1.
		\]
		由 $x\geqslant0,\ y\geqslant0,\ x+y\leqslant1$，
		以及两曲面 $z=1$、$z=1/(x+y)$，区域可写成
		\[
		0<s\leqslant1,\qquad 0\leqslant t\leqslant s,\qquad
		1\leqslant z\leqslant\frac1s.
		\]
		因此
		\[
		\begin{aligned}
		V
		&=\int_0^1\int_0^s\int_1^{1/s}dz\,dt\,ds\\
		&=\int_0^1s\left(\frac1s-1\right)ds
		=\int_0^1(1-s)\,ds
		=\boxed{\frac12}.
		\end{aligned}
		\]
		又因为在 $\Omega$ 内 $x^2+y^2+z^2\geqslant z^2$，所以
		\[
		\begin{aligned}
		\iiint_\Omega\frac{dx\,dy\,dz}{x^2+y^2+z^2}
		&\leqslant
		\int_0^1\int_0^s\int_1^{1/s}\frac{dz\,dt\,ds}{z^2}\\
		&=\int_0^1\int_0^s(1-s)\,dt\,ds\\
		&=\int_0^1s(1-s)\,ds
		=\frac16.
		\end{aligned}
		\]
		而 $\frac16<\frac14=\frac V2$，故题中不等式成立。
	\end{enumerate}
}


	
\end{enumerate}
